\documentclass[../../../include/open-logic-section]{subfiles}

\begin{document}

\olfileid{sth}{choice}{tarskiscott}
\olsection{Cara Tarski--Scott}

Ing \olref[cardinals][cardsasords]{defcardinalasordinal}, kita
ndhéfinisikaké kardinal minangka ordinal. Kanggo nindakaké iki, kita
nganggep Aksioma Pengurutan Becik. Alesané mung amarga cara iki
lumrah digunakaké ing disiplin iki.

Mula sadurungé ngrembug pitakonan filsafat babagan Pengurutan Becik,
perlu ditegesaké yèn kita \emph{bisa} ninggalaké cara kang lumrah mau
lan ngrembakakaké téori kardinal \emph{tanpa} nganggep Pengurutan
Becik. Kita isih bisa migunakaké dhéfinisi $\cardeq{A}{B}$,
$\cardle{A}{B}$ lan $\cardless{A}{B}$ kaya kang katon ing
\olref[sfr][siz][]{chap}. Kita mung mbutuhaké pangertèn anyar babagan
\emph{kardinal}.

Pamikiran naif yaiku nyoba ndhéfinisikaké kardinalitas $A$ kaya iki:
\[
	\Setabs{x}{\cardeq{A}{x}}.
\]
Panjenengan bisa mbandhingaké iki karo dhéfinisi Frege kanggo
$\# x Fx$, kang wis digambaraké ing pungkasan
\olref[cardinals][hp]{sec}. Amarga alasan kang wis kita tunjukaké ing
kana, dhéfinisi iki gagal kanggo himpunan kang ora kosong. Saben himpunan
singleton ekuinumerik karo $\{\emptyset\}$. Nanging singleton anyar
kawangun ing saben tahap penerus hirarki: cukup gatekna singleton
saka tahap sadurungé. Luwih umum, kanggo saben himpunan kang ora kosong,
prodhuk Kartesiusé karo singleton kang unsuré ana ing tahap saya dhuwur
tetep ekuinumerik karo himpunan mau, nanging pangkaté bisa saya dhuwur tanpa
wates. Mula $\Setabs{x}{\cardeq{A}{x}}$ ora dadi himpunan kanggo
himpunan kang ora kosong, amarga ora nduwèni wates pangkat. Yèn himpunan asal kosong,
kumpulan ekuinumerik mau pancèn dadi himpunan.

Kanggo ngatasi prakara iki, kita migunakaké cara saka Tarski lan Scott:\footnote{Élinga: saben rumus bisa nduwèni paramèter, kajaba yèn ana pratelan kang kanthi gamblang nglarang.}

\begin{defn}[Tarski--Scott]
Kanggo sembarang rumus $\phi(x)$, $[ x : \phi(x)] $ yaiku himpunan
kabeh $x$ kang pangkaté paling cilik ing antarané kang nyukupi
$\phi(x)$, utawa $\emptyset$ yèn ora ana kang nyukupi $\phi$.
\end{defn}

Kita kudu mriksa yèn dhéfinisi iki sah. Ing $\ZF$,
\olref[spine][foundation]{zfentailsregularity} njamin anané
$\setrank{x}$ kanggo saben $x$. Yèn ana obyèk kang nyukupi
$\phi$, kita bisa njupuk $\alpha$ minangka pangkat paling cilik
kang nyukupi $(\exists x\subseteq V_\alpha)\phi(x)$, yaiku
$(\forall \beta \in \alpha)(\forall x \subseteq V_\beta)\lnot \phi(x)$.
Banjur kanthi Pamisahan kita dhéfinisikaké $[x : \phi(x)]$ minangka
$\Setabs{x \in V_{\alpha+1}}{\phi(x)}$.

Sawisé mbeneraké cara Tarski--Scott, kita saiki bisa migunakaké
kanggo ndhéfinisikaké sawijining pangertèn kardinalitas:

\begin{defn}
Kardinalitas \textsc{ts} saka $A$ yaiku $\text{tsc}(A) = [x :
\cardeq{A}{x}]$.
\end{defn}

Dhéfinisi kardinal \textsc{ts} ora migunakaké Pengurutan Becik.
Nanging sanadyan tanpa aksioma iku, kita bisa nuduhaké yèn
\emph{kardinal \textsc{ts}} nduwèni sipat kang padha karo
\emph{kardinal} kaya kang didhéfinisikaké ing
\olref[cardinals][cardsasords]{defcardinalasordinal}.
Upamané, yèn \olref[cardinals][cardsasords]{lem:CardinalsBehaveRight}
lan \olref[card-arithmetic][opps]{lem:SizePowerset2Exp} kita
pratelakaké maneh nganggo kardinal \textsc{ts}, buktiné tetep lumaku
ing $\ZF$ tanpa nganggep Pengurutan Becik.

Nalika lagi ngrembug iki, perlu dielingaké yèn kita uga bisa
ngrembakakaké téori ordinal nganggo cara Tarski--Scott. Yèn
$\tuple{A, <}$ sawijining pengurutan becik, tetepna
$\text{tso}(A, <) = [\tuple{X, R} :
\ordeq{\tuple{A, <}}{\tuple{X, R}}]$. Katrangan luwih jembar
babagan cara iki kanggo kardinal lan ordinal bisa diwaca ing
\citet[bab~9--12]{Potter2004}.

\end{document}
